\section{Discusión}\label{sec:discusion}

\subsection{Qué mejoró realmente}
El modelo definitivo no se sostiene por añadir capas sin control, sino por la validez del experimento: corte cronológico, búsqueda multisemilla, regularización, comparadores bajo la misma política, calibración OOF, hashes y pruebas de fuga. La discriminación es moderada (PR-AUC 0.2249 frente a prevalencia 0.0995; ROC-AUC 0.7482), por lo que existe señal, pero no una separación fuerte.

\subsection{Red neuronal frente a modelos tabulares}
La MLP alcanza el mejor PR-AUC y ROC-AUC nominal de la comparación fija, mientras la regresión logística obtiene mejor F1 y \emph{recall} al umbral seleccionado. Esto es compatible con datasets tabulares medianos: la no linealidad puede mejorar el orden sin dominar todos los puntos operativos. Sin una prueba pareada predeclarada no corresponde afirmar superioridad estadística. El valor académico está en demostrar un experimento neuronal correcto y auditable, no en ocultar que un modelo lineal sigue siendo competitivo.

\subsection{Calibración no es discriminación}
Platt reduce Brier de 0.1817 a 0.0831 y ECE de 0.2623 a 0.0085, pero PR-AUC y ROC-AUC permanecen iguales. Esta distinción es central: calibrar hace interpretable la escala probabilística; no crea información ni corrige errores de ranking. El F1 calibrado ligeramente menor tampoco es contradicción: el umbral 0.20 proviene exclusivamente de validación OOF y no fue reajustado a 2024--2025.

\subsection{Errores y utilidad real}
Con 88 verdaderos positivos, 134 falsos negativos y 285 falsos positivos, la app no debe automatizar recursos, sanciones o responsabilidades. Su uso defendible es ordenar revisiones posteriores a la notificación y comunicar incertidumbre. Los casos cercanos al umbral evidencian que una etiqueta binaria oculta información; por eso la interfaz muestra probabilidad, umbral y referencias con intervalos.

\subsection{Sesgos y variables ausentes}
El registro ONSV proviene de partes administrativos, 2025 es preliminar y puede existir subregistro o cambio de codificación territorial. Además faltan ocupantes, velocidad estimada, tipo y antigüedad del vehículo, uso de cinturón/casco y exposición vial. Estas variables tienen relación mecánica directa con la multifatalidad y probablemente fijan el techo actual. Incorporarlas exige una fuente enlazable y disponible al momento de uso; no deben inventarse ni imputarse desde el desenlace.

\subsection{Alcance posterior al hecho}
La clase del evento y parte de la escena se conocen después de notificar el siniestro. En consecuencia, el modelo no predice qué vía tendrá un accidente, ni la probabilidad de fatalidad de un viaje. También queda prohibido extrapolar resultados a periodos futuros sin validación temporal nueva. Cuando exista un nuevo periodo realmente no observado, debe evaluarse una sola vez con el bundle actual antes de considerar cualquier actualización.

\subsection{Lectura correcta de SHAP}
Los grupos clase, red vial, tipo de vía$\times$zona y zona concentran la importancia global. Esto describe cómo el modelo distribuye su score en 512 casos de validación respecto de 256 casos de fondo. Los signos medios pueden ocultar categorías con direcciones opuestas; no son efectos causales ni recomendaciones de intervención. La ausencia deliberada de explicaciones locales evita presentar narrativas individuales no validadas.
